\begin{textsample}{NEG Dim 1 – vem\_ed\_21 – Score 4.00 – vem\_ed\_21/t108.txt}  \label{ex:f1_neg_013}
continuação Apoio ao matchmaking entre parceiros, treinamento docente, orientação e suporte, reconhecimento de professores e alunos estão entre as recomendações para institucionalizar as iniciativas COIL.
A segunda comunicação de Ana Carolina foi"Regional campuses in the lead: Institutionalizing COIL across multi-campus \textbf{public} systems"(Campi regionais na liderança: institucionalizando COIL em sistemas públicos multicampi), com Dan Nolan (University of Minnesota System / EUA) e Natalia Dyba (University of Washington Bothell / EUA).
Os autores relataram os esforços de institucionalização de \textbf{campi} regionais em sistemas públicos no Brasil e nos Estados Unidos.
Destacaram que a organização dos Intercâmbios Virtuais nesse contexto requer alinhamento com \textbf{metas} interinstitucionais mais amplas e uma noção dos pontos fortes, da preparação e da capacidade do corpo docente, da equipe e da liderança.

% matched lemmas: campi, meta, public
\end{textsample}
